\scanpage{117}
\mspage{1}
% archive sheet number: 116
Carrés $h$-cartésiens et $h$-cocartésiens

\deleted{Structures \unsure{quadrangulées|quadrangulaires}}

1. \uline{Axiomes} de \uline{structures} \uline{\unsure{quadrangulées|quadrangulaires}} (une tentative avortée)

% source lines 1-7
Soit : $(M,W)$ une catégorie de modèles admettant une théorie de la fibre homotopique. Alors on a dans $\mathrm{Ho}_W\unsure{\mathcal C|\mathscr H|\mathcal E}=MW^{-1}$ une notion de \uline{carré} \uline{$W$-cartésien} (ou $W$-homotopiquement cartésien), par quoi on entend un diagramme carré
\diagram{\begin{tikzcd}[ampersand replacement=\&]
X_1 \arrow[d] \& X_3 \arrow[l] \arrow[d] \\
X_0 \& X_2 \arrow[l]
\end{tikzcd}}
dans $\mathrm{Ho}_W$ provenant, à iso. près, d'un carré $W$-cartésien dans $M$ - donc \deleted{\unsure{par exemple}} isomorphe, dans $\underline{\mathrm{Hom}}(\unsure{\mathbb{\Delta}|\Delta}^1\times\mathbb{\Delta}^1,\unsure{\mathrm{HO}|\mathrm{Ho}}_W)$, à un carré \deleted{\unsure{fibré}} \inserted{\uline{cartésien}}
\deleted{\unsure{les flèches}}
\diagram{\begin{tikzcd}[ampersand replacement=\&, column sep=large]
\deleted{\unsure{X_1}} \arrow[d, draw=none, "{\deleted{\unsure{\downarrow}}}" description] \& Y_1 \arrow[d,"p_1"'] \& Y_3(\simeq Y_1\times_{Y_0}Y_2) \arrow[l] \arrow[d] \\
\deleted{\unsure{X_0}} \& Y_0 \& Y_2 \arrow[l,"p_1"]
\end{tikzcd}}
, avec $p_1,p_2\in\mathrm{Fib}_W$. (En fait, il suffit de $p_1$ ou $p_2\in\mathrm{Fib}_W$ pour assurer le caractère $W$-cartésien du carré cartésien envisagé.) Je voudrais axiomatiser
% margin emphasis beside the following line
\margin{!}
cette situation, \unsure{partons|partant} \unsure{en tentant} de faire la liste des propriétés \unsure{formelles|éventuelles} essentielles de cette notion, et de la notion duale de carré $W$-cocartésien (ou $W$-homotopiquement cocartésien), dans le cas où $(M,W)$ admet une théorie de la cofibre homotopique. L'\unsure{espèce|esquisse} de structure, \unsure{serait|savoir|voire} celle des ``catégories quadrangulées'' (à gauche, à \unsure{droit|droite}, bilatères) est une variante, dans un cas
